On a Dobot CR5, I learned a controller can appear faster yet become less reliable. Non-blocking execution at higher command frequency let each policy inference observe an intermediate arm state before the preceding action chunk finished. The resulting corrections drove the robot back and forth. I therefore used blocking, short-horizon ServoP action chunks, accepting lower apparent frequency to preserve state-action consistency and trajectory quality. This sharpened my guiding question: how can learning, control, timing, sensing, and safety form one reliable robot system?

My preparation combines complementary disciplines by design. I completed a B.Eng. in Mechanical Design, Manufacturing and Automation at Harbin Institute of Technology, Weihai, in June 2025. That degree and robotics projects grounded me in kinematics, robot modeling, and physical constraints. I am now pursuing a second bachelor's degree in Computer Science and Technology at Harbin Institute of Technology, Shenzhen, expected in June 2027. This is not a departure from mechanical engineering, but a deliberate expansion from modeling machines to building the algorithms, data pipelines, and interfaces through which they perceive, learn, and act. Together, these perspectives let me examine embodied systems from physical and computational sides.

That combination became concrete in a CR5 data-collection and inference loop using an Orbbec Astra2 RGB-D camera, an electric gripper, and ROS2. I anchored HDF5 synchronization to RGB timestamps, then aligned robot states, target actions, and gripper feedback through binary search and linear interpolation. Using Fast DDS shared memory and RGB-only capture, I reduced recorded frame drops from 27.15\% to the 2.57\%-4.74\% range. Rather than equate throughput with quality, I checked schema and count consistency, motion jumps, workspace bounds, start-pose distance, and slow playback before using or replaying an episode. The pipeline mapped HDF5 images, seven-dimensional robot-plus-gripper states and actions, timestamps, and episode, frame, and task indexes into LeRobot v2.0 artifacts. For deployment, I adapted OpenPI with matching mappings and WebSocket serving; its client built observations from RGB, actual robot state, gripper feedback, and a task prompt, then executed arm and gripper targets through ServoP and Modbus. This work clarified that data engineering is part of learning itself: clock selection, alignment, and state-action definitions determine the examples a policy receives, while validation determines which episodes should be used or replayed.

That question led me to action representation. As a Research Assistant at iLearn, I am second author of Hermite Action Tokenizer for VLA, an unsubmitted working paper, and I am responsible for its discrete autoregressive tokenization branch. Working on this branch made representation a control question: what trajectory structure should tokenization preserve, and how should gripper decisions be encoded separately? The method separates seven-dimensional action chunks into six-dimensional end-effector trajectories and gripper states, with shared breakpoints removing endpoint redundancy and quantization seams. In an offline codec benchmark averaging four LIBERO subsets with 8 by 7 action chunks, Hermite's Endpoint MSE was approximately 99.7\% lower than FAST and 82.6\% lower than BEAST. I keep that result within its protocol: it measures offline codec fidelity, not online rollout success or physical-robot error. The research showed that action encoding constrains what a model can express; CR5 deployment showed that execution schedules determine which state the model observes and what its command controls. Reliable embodied intelligence therefore depends on both representation semantics and execution semantics, joined through synchronized data.

During the Xbotics Embodied AI Community Internship, I migrated Isaac Lab's ANYmal-C navigation task into a MotrixLab NumPy environment. My judgment was to treat migration as interface preservation, not transcription. I refactored reset and step flows, command sampling, observations, rewards, and termination, corrected MuJoCo heightfield frames, added spawn and goal constraints, and used reward and termination regression checks. Earlier, in the HIT Weihai HERO RoboMaster Team's Vision Group, I helped modularize image acquisition, vision processing, and communication synchronization as ROS2 nodes and worked on host-MCU timing. These experiences reinforced that coordinate, interface, and timing assumptions must be explicit for integration.

My immediate goal is robotics or embodied-AI R\&D engineering, connecting learning methods to sensing, control, and deployment. I want to strengthen my ability to design and evaluate real-hardware systems, where representation, timing, and physical constraints must be reconciled. After gaining stronger real-system experience, I may pursue doctoral study; it remains an option rather than a predetermined destination. I therefore seek graduate training whose curriculum and applied opportunities can close gaps between my integration experience and the robot systems I want to build.
